\section{Recovery from four values}\label{sec:recovery}

The following result separates the local construction from the
degree automorphism to which it will be applied.

\begin{proposition}[Recovery from four values]\label{prop:recovery}
Let $F:\R\to2^\N$ be Borel with countable fibers, and suppose
\[
F(x+y)\leT F(x)\oplus F(y)\qquad(x,y\in\R).
\]
For every irrational $t\in(0,1)$ there is a comeager set
$G_t\subseteq\R^2$ such that for each $(u,v)\in G_t$ there
is a rational $\delta\in(0,1)$ for which
\begin{equation}\label{eq:four-values}
C(t)\leT
F(u)\oplus F(u-v)\oplus
F(v-u+\delta t)\oplus F(v-u+\delta(t-1)).
\end{equation}
\end{proposition}

\begin{proof}
The two auxiliary reals allow us to advance from $u+s$ to
$u+s+a$ through two additions:
\[
(u+s)+(v-u+a)=v+s+a,\qquad
(v+s+a)+(u-v)=u+s+a.
\]
If the programs witnessing these additions can be kept fixed,
then the oracles $F(v-u+a)$ and $F(u-v)$ turn an oracle for
$F(u+s)$ into one for $F(u+s+a)$. We will use only the two
increments $\delta t$ and $\delta(t-1)$. The initial oracle
$F(u)$, the fixed return oracle $F(u-v)$, and the two increment
oracles are exactly the four values in \eqref{eq:four-values}.
For each pair $(u,v)$ in a comeager set, we first arrange two
programs that remain fixed along the recurrence. A bit of
$F(u+s)$ will then select the increment so that successive
states stay bounded; the frequencies of the choices will
recover $t$.

\medskip
\noindent\emph{Two local addition programs.}
For each $(x,y)$, let $e(x,y)$ be the least index satisfying
\[
\Phi_{e(x,y)}^{F(x)\oplus F(y)}=F(x+y).
\]
Such an index exists by hypothesis. For a fixed index, the
condition that every output bit is computed correctly is
arithmetical in the three oracle sets and is therefore Borel
in $(x,y)$. The set on which this is the least successful
index is Borel as well. Thus $e:\R^2\to\N$ is Borel.

By Lemma~\ref{lem:continuity}, choose comeager sets
$D\subseteq\R$ and $E\subseteq\R^2$ such that $F|D$ and
$e|E$ are continuous; the codomain of $e$ is discrete.
Fix the given $t$ and put $H=\Q+\Q t$. Let $G_t$ consist
of the pairs $(u,v)$ satisfying
\begin{equation}\label{eq:all-shifts}
u+s\in D,\qquad
(u+s,v-u+a)\in E,\qquad
(v+s,u-v)\in E
\quad\text{for all }s,a\in H.
\end{equation}
The group $H$ is countable. For each fixed pair of shifts,
the last two maps of $(u,v)$ in \eqref{eq:all-shifts}
are affine homeomorphisms of $\R^2$. The first condition
is a pullback under a translated coordinate projection;
the inverse image of a meager exceptional set is meager.
Consequently $G_t$ is comeager.

Fix any $(u,v)\in G_t$, and set
\[
p=e(u,v-u),\qquad q=e(v,u-v).
\]
Both centers belong to $E$. Relative continuity of $e|E$
gives a number $0<\rho<1$ such that
\begin{align}
e(u+s,v-u+a)&=p
&& (s,a\in H,\ |s|,|a|<\rho),\label{eq:p}\\
e(v+s,u-v)&=q
&& (s\in H,\ |s|<\rho).\label{eq:q}
\end{align}
The membership in $E$ required for these equalities is
supplied by \eqref{eq:all-shifts}. Thus $p$ and $q$ carry out
the two additions above whenever $s,a\in H$ and
$|s|,|a|,|s+a|<\rho$. The last bound is needed because the
second addition starts at the new offset $s+a$. We now choose
the two increments and a rule for selecting between them so
that the old and new offsets always stay in this range.

\medskip
\noindent\emph{A bounded recurrence.}
Choose $-\rho<b_-<0$ with $u+b_-\in D$. The interval
$(0,\rho)$ contains uncountably many $b$ with $u+b\in D$,
and the fiber of $F(u+b_-)$ is countable. We may therefore
choose $0<b_+<\rho$ with $u+b_+\in D$ and
$F(u+b_-)\ne F(u+b_+)$. Fix a bit $j$ where these two
sets differ, and put $h=F(u+b_+)(j)$.

Continuity of $F|D$ makes this bit constant on sufficiently
small relative neighborhoods of the two endpoint translates.
Choose a positive rational $\delta<\rho$ small enough that
\begin{align}
F(u+s)(j)&\ne h
&&\bigl(s\in H\cap[b_-,b_-+\delta]\bigr),\label{eq:left}\\
F(u+s)(j)&=h
&&\bigl(s\in H\cap[b_+-\delta,b_+]\bigr).\label{eq:right}
\end{align}
We may also require $3\delta<b_+-b_-$. Every tested
translate belongs to $D$ by \eqref{eq:all-shifts}; the
endpoints themselves need not lie in $H$.

Define $s_0=0$ and, recursively,
\begin{equation}\label{eq:recurrence}
r_n=
\begin{cases}
1,&F(u+s_n)(j)=h,\\
0,&F(u+s_n)(j)\ne h,
\end{cases}
\qquad
s_{n+1}=s_n+\delta(t-r_n).
\end{equation}
We claim that
\begin{equation}\label{eq:invariant}
s_n\in H\cap[b_-,b_+]\qquad(n\ge0).
\end{equation}
The initial state has this property. If $r_n=0$, then
\eqref{eq:right} forces $s_n<b_+-\delta$; the increment
$\delta t$ lies in $(0,\delta)$, so $s_{n+1}$ remains
in the interval. If $r_n=1$, then \eqref{eq:left} forces
$s_n>b_-+\delta$; its negative increment has magnitude
$\delta(1-t)<\delta$, with the same conclusion.
Rationality of $\delta$ ensures $s_{n+1}\in H$.
This proves \eqref{eq:invariant}, and in particular
$|s_n|<\rho$.
Figure~\ref{fig:bounded-recurrence} isolates the only feature of the
tested bit needed for this confinement: it points the next step
inward near either endpoint. No restriction on its intervening
behavior is needed.
\begin{figure}[tb]
\centering
\begin{tikzpicture}[x=1cm,y=1cm,>=Stealth,font=\small]
  \fill[blue!12] (0,0) rectangle (1.2,.55);
  \fill[gray!7] (1.2,0) rectangle (8.8,.55);
  \fill[orange!18] (8.8,0) rectangle (10,.55);
  \draw[thick] (0,0) -- (10,0);
  \foreach \x in {0,1.2,8.8,10} {\draw (\x,-.08) -- (\x,.62);}
  \node[below] at (0,-.09) {$b_-$};
  \node[below] at (1.2,-.09) {$b_-+\delta$};
  \node[below] at (8.8,-.09) {$b_+-\delta$};
  \node[below] at (10,-.09) {$b_+$};
  \node at (.6,.3) {$r_n=0$};
  \node at (9.4,.3) {$r_n=1$};
  \node at (5,.3) {either choice may occur};
  \draw[->,thick,blue!65!black] (.12,.9) -- (1.1,.9);
  \node[above,blue!65!black] at (.61,.97) {$+\delta t$};
  \draw[->,thick,orange!70!black] (9.88,.9) -- (8.9,.9);
  \node[above,orange!70!black] at (9.39,.97) {$-\delta(1-t)$};
\end{tikzpicture}
\caption{The bounded recurrence, schematically. The tested bit forces
an inward step in each endpoint strip; the bit need not be monotone
between them. Each step has length less than $\delta$, so no state
can leave $[b_-,b_+]$. The oracle computation produces the choices
$r_n$, not the numerical states $s_n$.}
\label{fig:bounded-recurrence}
\end{figure}


\medskip
\noindent\emph{Computing the binary choices.}
Let $T$ denote the finite join on the right of
\eqref{eq:four-values}. We show that $(r_n)_{n\ge0}$ is
computable from $T$ by maintaining, relative to $T$, an
oracle-program index $P_n$ satisfying
\[
\Phi_{P_n}^{T}=F(u+s_n).
\]
Initially $P_0$ projects to the first component $F(u)$.
Given $P_n$, read $\Phi_{P_n}^{T}(j)$ and compare it with
the fixed bit $h$ to obtain $r_n$. Select the third component
of $T$ if $r_n=0$ and the fourth if $r_n=1$, namely
\[
F(v-u+\delta(t-r_n)).
\]

Set $a_n=\delta(t-r_n)$. By \eqref{eq:invariant},
$s_n,a_n\in H$ and $|s_n|,|a_n|<\rho$. Equation~\eqref{eq:p}
therefore shows that program $p$, with its two oracle
components simulated by $\Phi_{P_n}^{T}$ and the selected
tuple component, computes
\[
F(v+s_n+a_n)=F(v+s_{n+1}).
\]
Since $|s_{n+1}|<\rho$, Equation~\eqref{eq:q} then shows
that program $q$, with this output and the fixed component
$F(u-v)$ as its two oracles, computes $F(u+s_{n+1})$.
Effective substitution of oracle subroutines produces
the next index $P_{n+1}$.

This is an actual computation relative to the fixed tuple.
The indices are constructed using the already computed
branch bits. Inductively each substituted subroutine is
total on its actual oracle, and both outer programs are
valid reductions at the indicated argument pairs. Every
bit computation has finite nesting depth through earlier
indices and terminates. No bound on its running time is
needed. To compute $r_n$, perform the finite construction
through stage $n$.

The program uses only the fixed integers $p,q,j,h$, the
rational $\delta$, and the four tuple components. It does
not compute the numerical states $s_n$, or use
$t,u,v,\rho,b_-,b_+$ as additional real oracles. The finite
program constants may depend on $t,u,v$, as allowed in the
pointwise reduction \eqref{eq:four-values}.

\medskip
\noindent\emph{Recovering the parameter.}
Summing \eqref{eq:recurrence} gives
\[
s_N=\delta\left(Nt-\sum_{n<N}r_n\right).
\]
Because $|s_N|<\rho<1$, the $T$-computable rational numbers
$A_N=N^{-1}\sum_{n<N}r_n$ satisfy
\begin{equation}\label{eq:error}
|t-A_N|<\frac{1}{N\delta}\qquad(N\ge1).
\end{equation}
For a rational query $a$, search for an $N$ such that
$a<A_N-1/(N\delta)$ or $a>A_N+1/(N\delta)$, and return
the corresponding cut bit. Irrationality of $t$ guarantees
termination. The rational $\delta$ makes the error bound
effective, so this computes $C(t)$ from $T$ and proves
\eqref{eq:four-values}.
\end{proof}

\begin{remark}\label{rem:nonuniform}
The parameter $\delta$ and the reduction program in
Proposition~\ref{prop:recovery} need not be selected
effectively from $(u,v)$. The conclusion is that the
reduction exists for every pair in a comeager set.
The category avoidance argument allows precisely this
form of nonuniformity.
\end{remark}
